% Fragmento científico exclusivo de 01_electron.tex, líneas 120--197.
% El prefijo 1--118 coincide byte a byte con la ontología principal ya
% publicada; este fragmento conserva únicamente la formulación cuántica
% adicional sin repetir definiciones idénticas.
\section{Ontología cuántica de la partícula}
\Needspace{7\baselineskip}
\subsection{Estado, historia y realización}
Un estado genealógico es una sección compatible del sistema de extensiones del TPK. Una historia observable es la proyección cronológica de esa sección sobre un lector. Una realización física es la representación de la historia en una recta dimensional, un espacio de Hilbert, un álgebra de campos o un codominio topológico. Estas tres nociones no se intercambian. El estado porta la identidad; la historia porta la evidencia; la realización porta el observable.

Sea \(\mathcal X_{\rm TPK}^{\rm enr}\) el espacio de estados enriquecidos y \(p\) el lector observable. Una partícula es una clase \([X]\) de secciones persistentes tal que

\[
p(X_{K+1})|_K=p(X_K),
\]
la ley de transición conserva las relaciones del registro genealógico y la prolongación no termina en un horizonte de obstrucción. La ruta

\[
\gamma_X=(p(X_0),p(X_1),p(X_2),\ldots)
\]
es la historia visible de la partícula. No es una trayectoria espacial por definición. Una realización métrica posterior puede convertirla en curva, órbita o propagador cuando construye el entrelazador correspondiente.

La identidad física se conserva bajo refinamiento cuando los cuadrados

\[
\begin{CD}
X_{K+1} @>{F_{K+1}}>> X'_{K+1}\\
@V{r_{K+1,K}}VV @VV{r'_{K+1,K}}V\\
X_K @>>{F_K}> X'_K
\end{CD}
\]
conmutan. La masa, la carga y el espín son lectores diferentes de ese mismo antecedente. La igualdad de uno no obliga a la igualdad de los demás.

\Needspace{7\baselineskip}
\subsection{Representación cuántica de la partícula como sección local persistente del estado genealógico}
Una partícula HMT--MD es una sección persistente de una extensión enriquecida que satisface cuatro condiciones:

\begin{enumerate}
\item posee una ruta observable compatible con la monodromía nonádica;
\item conserva un registro genealógico bajo refinamiento;
\item admite una representación de orientación, carga, espín y estadística;
\item presenta una salida espectral estable en el codominio físico declarado.
\end{enumerate}
La estabilidad no significa inmovilidad. Significa que los cambios de fase, hoja y memoria pertenecen a una misma clase genealógica. La partícula puede cerrar su fase y no su estado, cambiar de hoja y conservar su masa, o atravesar varias celdas y mantener la misma identidad de extensión.

Una especie es una familia de partículas que comparte el clasificador interno

\[
\mathfrak s=
(\mathfrak q,\mathfrak j,\mathfrak c,\mathfrak f,\mathfrak g,
\mathfrak r,\mathfrak o),
\]
donde \(\mathfrak q\) es la representación de carga, \(\mathfrak j\) la holonomía de espín, \(\mathfrak c\) la órbita de color, \(\mathfrak f\) el sabor, \(\mathfrak g\) la generación, \(\mathfrak r\) la familia de ruta y \(\mathfrak o\) el operador de masa. Una especie no central puede corresponder a una multisección completa; exigirle un único punto APP antes de declarar el acoplamiento físico constituye un error de tipo.

\Needspace{7\baselineskip}
\subsection{Conjugación cuántica de la antipartícula mediante inversión de hoja de una sección persistente}
La involución dodecafásica es

\[
C_M(\tau)=-\operatorname{rev}(\tau).
\]
Actúa sobre la historia orientada y se prolonga al registro genealógico. La carga y la orientación lineal son impares; la memoria cuadrática es par. Si \(q\) es el lector de carga,

\[
q(C_MX)=-q(X).
\]
La igualdad de masas exige además covariancia operatoria. Si \(K\) realiza la conjugación en el espacio físico y

\[
K\widehat M_XK^{-1}=\widehat M_{\bar X},
\]
transportando también los canales y la prescripción causal cuando \(K\) es antiunitario, entonces \(m_{\bar X}=m_X\). La paridad del registro por sí sola no sustituye esta premisa.

Partícula y antipartícula son dos secciones orientadas de una misma banda genealógica. En la realización de Möbius, una vuelta intercambia la hoja y dos vueltas restauran la orientación. La conclusión no procede de dibujar una cinta torcida: exige la involución, el levantamiento a la doble cubierta, el lazo de retorno y la holonomía de signo.

Para el registro electrónico \(\tau_e=+++---+++---\),

\[
C_M\tau_e=\tau_e.
\]
La palabra central es fija, mientras la representación de carga distingue electrón y positrón. Así se obtiene una misma masa sin duplicar la ruta temporal. El cierre a 216 pasos no significa 108 pasos de partícula y 108 de antipartícula: significa restauración del levantamiento cinemático orientado.

\Needspace{7\baselineskip}
